\section{Repaso: las limitaciones del gradiente de políticas}

El gradiente de políticas usa el reward-to-go muestreado para estimar qué tan buena o mala es cada acción, pero esta estimación tiene \textbf{varianza muy alta}.

\begin{warningbox}{Los dos problemas centrales del gradiente de políticas}
\begin{enumerate}
    \item \textbf{Aprovechamiento insuficiente de los datos}: bajo recompensas dispersas, la recompensa de la mayoría de las trayectorias es cero, y la información del gradiente es escasa.
    \item \textbf{Varianza alta}: el reward-to-go muestreado es una estimación de una sola muestra, con mucho ruido.
\end{enumerate}
\end{warningbox}

Idea central de la mejora: ¿es posible \textbf{aprender} qué es bueno y qué es malo, en lugar de depender únicamente de la estimación muestreada?

\subsection{Resumen de la sección}

El principal enemigo del gradiente de políticas es la varianza alta y el aprovechamiento ineficiente de las trayectorias; construir un critic aprendido para estimar la función de ventaja puede mejorar considerablemente este fenómeno, y sienta las bases para la aparición de los sistemas Actor-Critic.

